% SOURCE: OCW L11--12; Johnson 2008 matrix multiplication performance and cache-oblivious experiments, ideal cache terminology.
\originalchapter{计算成本与数据移动}\label{ch:performance}
\originalsection{操作数为什么不能预测运行时间}
两个矩阵乘法程序都按 $c_{ij}=\sum_k a_{ik}b_{kj}$ 计算, 都有约 $2n^3$ 次实数加乘, 实际速度却可能相差很大. 原课程的实验是在 2.66GHz Intel Core 2 Duo、gcc 4.1.2 与 ATLAS 3.6.0 的历史平台上进行, 展示缓存容量、访问顺序和寄存器复用的影响. 这些曲线不能被当作今日硬件的测试结果, 但它们提出的问题仍适用于现代算法设计: 数据从哪里来, 会被使用几次, 下一次使用前是否已经被替换?

算 $q(q^*a)$ 时, 先算长度 $m$ 的点积再作向量缩放, 成本为 $O(m)$. 若先形成 $qq^*$ 再乘 $a$, 要存储 $m^2$ 个元素并作 $O(m^2)$ 运算. 矩阵乘法在精确算术中结合, 但括号的位置既改变成本, 也可能改变舍入路径. 数学表达式必须被转化为恰当的计算图.

\begin{definition}
在理想二级存储模型中, 快速存储可容纳 $Z$ 个数, 主存容量充足. 需要的数据已在快速存储中称命中, 否则称缺失, 需要把数据搬入并可能替换旧数据. 完全相联表示任意数据都可放到任意位置; 最优替换使用未来访问信息, 淘汰最久不会再用的数据. 缓存复杂度 $Q(n;Z)$ 计数据移动量或缺失次数, 必须先说明一个移动单位是一个数还是一个缓存行.
\end{definition}
现实缓存有有限相联度、多个层级和预取, 不能与理想模型逐一等同. 理想模型用于看清主要数量级, 不保证常数和每个矩阵尺寸下的性能. 时间局部性指一个值在短时间内反复用; 空间局部性指相邻地址在短时间内用到, 使一整条缓存行有效.

\originalsection{行列存储与访存方向}
行主序按每一行连续存储, 列主序按每一列连续存储. 设缓存行含 $B$ 个数. 顺序读 $n$ 个数约需 $n/B$ 次搬运; 若访问步长很大, 每次只用一条缓存行中的一个数, 可能接近 $n$ 次搬运. 运算次数相同, 带宽成本却差约 $B$ 倍.

Julia 的通常稠密数组采用列主序. 所以对矩阵逐元素遍历时, 内层循环宜让行指标改变, 沿一列连续走. C 中常见二维数组采用行主序, 适宜方向恰好相反. 语言和库的存储约定要实查, 不能凭矩阵的数学行列想象内存布局.

\begin{example}
一个 $n\times n$ 矩阵按列主序存储, 连续缓存行有 $B$ 个数. 比较逐列与逐行扫描全部元素的缺失次数, 假设步长和缓存容量使逐行访问不能保留有用的缓存行.
\end{example}
\begin{solution}
逐列扫描沿连续地址, 每条缓存行中的 $B$ 个数都被使用, 约为 $n^2/B$ 次. 逐行扫描每次跳过 $n$ 个数, 若缓存相联和容量使行间不能复用, 每个元素都可能触发一次搬入, 约为 $n^2$ 次. 这里的第二个数量级带有明确模型假设; 若整个矩阵能放入缓存, 或预取和跨行复用有效, 实测会不同.
\end{solution}
真实缓存的地址映射还可能造成某些特殊尺寸的冲突缺失. 因此时间曲线中的尖峰不一定是随机噪声. 测性能时应该记录尺寸、布局、线程、库与重复测量, 而不是只给一个最快时间.

\originalsection{分块矩阵乘法}
对 $n\times n$ 稠密矩阵, 把各矩阵分为 $b\times b$ 块. 以块指标 $I,J,K$ 写
\[
C_{IJ}\leftarrow C_{IJ}+A_{IK}B_{KJ}.
\]
若 $3b^2\le Z$, 三个活动块可同时进入快速存储. 这时每个读入的 $A$、$B$ 元素能服务一个块内的多个乘加, 不再只用一次就被替换.

\begin{proposition}\label{prop:blockcache}
在以一个数为移动单位的理想模型中, 取 $b$ 与 $\sqrt Z$ 同阶、并使三个活动块能装入快速存储. 经典分块矩阵乘法的数据移动量为
\[
O\left(n^2+\frac{n^3}{\sqrt Z}\right).
\]
若按含 $B$ 个数的连续缓存行计数, 且块布局与 tall-cache 等条件支持有效连续搬运, 相应量级为 $O(n^2/B+n^3/(B\sqrt Z))$.
\end{proposition}
\begin{proof}
忽略边缘不整块的常数, 共有 $(n/b)^3$ 次块乘法. 每次搬运至多常数个 $b^2$ 规模块, 总量为 $O(n^3/b)$. 初始读入与最终写出至少涉及 $O(n^2)$ 个元素. 取 $b=\Theta(\sqrt Z)$ 即得. 缓存行版本要求活动块的有效访问沿连续地址发生, 每次移动的一行能利用 $\Theta(B)$ 个数, 所以相同元素搬运界除以 $B$; 若这个布局条件失败, 不能直接这么除.
\end{proof}
这个算法需要知道适当块大小. 多级缓存需要多层分块; 处理器寄存器又构成更小的快速存储层. 高性能 BLAS 往往采用缓存块、数据打包和寄存器小核的结合, 不仅是给三重循环外面再加三层循环.

\originalsection{为什么平方根尺度自然出现}
$\sqrt Z$ 不是随意选的参数. 经典矩阵乘法的每个标量乘积 $a_{ik}b_{kj}$ 需要一对输入并更新输出位置 $(i,j)$. 把一次阶段中完成的乘积索引集合记为 $T\subset\{(i,k,j)\}$, 对应三个投影集合是用过的 $A$、$B$、$C$ 位置.

\begin{lemma}
若这些三个投影集合分别有 $a,b,c$ 个位置, 则 $|T|\le\sqrt{abc}$.
\end{lemma}
\begin{proof}
用 $0$--$1$ 矩阵 $M,N,P$ 分别指示三个位置集合. $T$ 的每个三元组都被这三个指标包含, 所以
\[
|T|\le\sum_{i,k,j}M_{ik}N_{kj}P_{ij}
=\inner{P}{MN}_F
\le\norm{P}_F\norm{MN}_F
\le\norm{P}_F\norm{M}_F\norm{N}_F=\sqrt{abc}.
\]
这里矩阵 Frobenius 内积的 Cauchy--Schwarz 不等式给第一界; $\norm{MN}_F\le\norm{M}_F\norm{N}_F$ 可逐列应用算子范数界与 $\norm{M}_2\le\norm{M}_F$ 得到.
\end{proof}
在经典不重复计算、显式维护各输出部分和的模型中, 把执行划成每段至多 $Z$ 次元素搬运. 一段初始快存有 $Z$ 个数, 搬运再引入至多 $Z$ 个, 所以输入位置的总规模为 $O(Z)$. 输出部分和可能从零新建, 不能只数搬入: 每个被更新的输出位置或者在段初已存在, 或在段中被搬入/搬出, 或在段末仍驻留快存. 由于不能丢弃待保留的部分和且不允许重计算, 三类数量均为 $O(Z)$, 输出位置也只有 $O(Z)$ 个. 因此三个投影位置集合的总规模为 $O(Z)$. 引理给该段最多 $O(Z^{3/2})$ 个不同乘积. 要完成 $n^3$ 个乘积, 需要 $\Omega(n^3/Z^{3/2})$ 段, 每完整段搬运 $Z$ 个数, 给出 $\Omega(n^3/\sqrt Z)$ 主项, 再加输入输出的 $\Omega(n^2)$.

这是经典矩阵乘法存储下界的主要结构解释. 若使用快速矩阵乘法、允许大量重计算, 或换成不同数据移动模型, 下界形式必须重新论证. 不能由此宣称任何可能的矩阵乘法算法都至少要 $n^3$ 次算术.

\originalsection{缓存无关的递归算法}
缓存无关算法不把 $Z$ 或 $B$ 写进程序. 对矩阵乘法, 把三个维度分成两半, 用分块恒等式递归计算, 直到足够小的基础问题. 对方阵, 一层产生八个半尺寸乘法, 总算术仍是 $\Theta(n^3)$.

当递归子问题的三个矩阵块达到能同时放入快存的尺寸 $b=\Theta(\sqrt Z)$ 时, 以下层级的访问都能在快存中完成. 该算法虽然不知道 $Z$, 递归树却自动经过这一尺度. 共有约 $(n/b)^3$ 个这种叶块, 每个搬运 $O(b^2)$ 个数, 因而也得到 $O(n^3/\sqrt Z+n^2)$ 元素搬运. 缓存行版本仍需恰当布局与容量条件.

对很不方的矩阵, 通常优先切最长维度, 避免产生极细长、不利于复用的块. 实际递归不应一直到 $1\times1$: 函数调用开销、寄存器利用与向量化可能压过理论优势. 原课程手册用较小块作为基例, 并指出即便消除 L1/L2 的明显容量断崖, 未优化的递归代码仍可能明显慢于专门 BLAS. 因而“缓存无关的渐近最优”不等于“任何尺寸的实际速度最快”.

\begin{center}
\begin{tikzpicture}[scale=.85,>=Stealth]
\draw (0,0) rectangle (2,2);\draw (1,0)--(1,2);\draw (0,1)--(2,1);
\node at (1,-.35) {$A$};
\draw (3,0) rectangle (5,2);\draw (4,0)--(4,2);\draw (3,1)--(5,1);
\node at (4,-.35) {$B$};
\draw (6,0) rectangle (8,2);\draw (7,0)--(7,2);\draw (6,1)--(8,1);
\node at (7,-.35) {$C$};
\node at (2.5,1) {$\times$};\node at (5.5,1) {$=$};
\node at (4,2.45) {$C_{11}=A_{11}B_{11}+A_{12}B_{21}$};
\end{tikzpicture}
\end{center}
图中的分块等式是精确矩阵恒等式. 不同求和树在浮点中可能给不同末位, 因此性能优化也应搭配误差与残差检查.

\originalsection{BLAS 层级与分解的块实现}
BLAS 一级操作是点积、向量加法等, 每个数据通常只参与少量运算; 二级操作是矩阵向量乘法与秩一更新; 三级操作是矩阵乘法和多右端三角求解, 能用大量数据复用提高算术强度. 算术强度是每搬运一个字节完成的操作数. 当强度低时, 运算往往受内存带宽限制; 当强度高且小核优化合理时, 更可能接近计算吞吐极限.

分块 LU 先分解一个窄面板, 再用三角求解和矩阵乘法更新尾块. 分块 Cholesky 对对角块分解, 解下方块, 再作对称秩更新. 分块 QR 把多个反射表示为 $I-VTV^*$, 利用矩阵乘法成批作用. 它们保留原有分解的数学目标, 但把大部分工作移到三级 BLAS. 若为重排计算改变了舍入次序, 应核查原算法的稳定性分析是否仍适用.

可比较的性能实验至少要先预热编译和库初始化, 固定线程设置, 用足够多次重复区分波动, 把输入生成和文件操作从计时中分离, 并检查输出正确性. 所谓某个程序“快十倍”应附具体环境和规模, 不能从历史课件曲线直接推给本次实现.

\begin{exercise}
若 $Z=3b^2$, 比较 $n\gg b$ 时逐元素矩阵乘法的 $O(n^3)$ 元素搬运与分块算法的 $O(n^3/b)$ 搬运. 当快存容量增加四倍时, 分块算法的主导搬运项如何改变?
\end{exercise}
\begin{solution}
$b$ 与 $\sqrt Z$ 同阶, 因而 $Z$ 增加四倍使合适的 $b$ 增加两倍, 主导 $n^3/b$ 项约减半. 这只是在同一理想模型中的主项比较, 不保证实际墙钟时间也恰好减半.
\end{solution}
\begin{remark}
本章整理 L11--12 的三份性能手册及逐周说明. 历史硬件、实验现象、寄存器复用和缓存无关观点来自 Johnson 手册; 分块上下界的中文推导和图为自编补充. 本册没有把原手册的旧性能曲线改标为新实测.
\end{remark}
